% Options for packages loaded elsewhere
% Options for packages loaded elsewhere
\PassOptionsToPackage{unicode}{hyperref}
\PassOptionsToPackage{hyphens}{url}
%
\documentclass[
  ignorenonframetext,
  aspectratio=169,
]{beamer}
\newif\ifbibliography
\usepackage{pgfpages}
\setbeamertemplate{caption}[numbered]
\setbeamertemplate{caption label separator}{: }
\setbeamercolor{caption name}{fg=normal text.fg}
\beamertemplatenavigationsymbolshorizontal
% remove section numbering
\setbeamertemplate{part page}{
  \centering
  \begin{beamercolorbox}[sep=16pt,center]{part title}
    \usebeamerfont{part title}\insertpart\par
  \end{beamercolorbox}
}
\setbeamertemplate{section page}{
  \centering
  \begin{beamercolorbox}[sep=12pt,center]{section title}
    \usebeamerfont{section title}\insertsection\par
  \end{beamercolorbox}
}
\setbeamertemplate{subsection page}{
  \centering
  \begin{beamercolorbox}[sep=8pt,center]{subsection title}
    \usebeamerfont{subsection title}\insertsubsection\par
  \end{beamercolorbox}
}
% Prevent slide breaks in the middle of a paragraph
\widowpenalties 1 10000
\raggedbottom
\AtBeginPart{
  \frame{\partpage}
}
\AtBeginSection{
  \ifbibliography
  \else
    \frame{\sectionpage}
  \fi
}
\AtBeginSubsection{
  \frame{\subsectionpage}
}
\usepackage{iftex}
\ifPDFTeX
  \usepackage[T1]{fontenc}
  \usepackage[utf8]{inputenc}
  \usepackage{textcomp} % provide euro and other symbols
\else % if luatex or xetex
  \usepackage{unicode-math} % this also loads fontspec
  \defaultfontfeatures{Scale=MatchLowercase}
  \defaultfontfeatures[\rmfamily]{Ligatures=TeX,Scale=1}
\fi
\usepackage{lmodern}

\usetheme[]{metropolis}
\usecolortheme[]{dolphin}
\usefonttheme[]{professionalfonts}
\ifPDFTeX\else
  % xetex/luatex font selection
\fi
% Use upquote if available, for straight quotes in verbatim environments
\IfFileExists{upquote.sty}{\usepackage{upquote}}{}
\IfFileExists{microtype.sty}{% use microtype if available
  \usepackage[]{microtype}
  \UseMicrotypeSet[protrusion]{basicmath} % disable protrusion for tt fonts
}{}
\makeatletter
\@ifundefined{KOMAClassName}{% if non-KOMA class
  \IfFileExists{parskip.sty}{%
    \usepackage{parskip}
  }{% else
    \setlength{\parindent}{0pt}
    \setlength{\parskip}{6pt plus 2pt minus 1pt}}
}{% if KOMA class
  \KOMAoptions{parskip=half}}
\makeatother

\usepackage{color}
\usepackage{fancyvrb}
\newcommand{\VerbBar}{|}
\newcommand{\VERB}{\Verb[commandchars=\\\{\}]}
\DefineVerbatimEnvironment{Highlighting}{Verbatim}{commandchars=\\\{\}}
% Add ',fontsize=\small' for more characters per line
\usepackage{framed}
\definecolor{shadecolor}{RGB}{241,243,245}
\newenvironment{Shaded}{\begin{snugshade}}{\end{snugshade}}
\newcommand{\AlertTok}[1]{\textcolor[rgb]{0.68,0.00,0.00}{#1}}
\newcommand{\AnnotationTok}[1]{\textcolor[rgb]{0.37,0.37,0.37}{#1}}
\newcommand{\AttributeTok}[1]{\textcolor[rgb]{0.40,0.45,0.13}{#1}}
\newcommand{\BaseNTok}[1]{\textcolor[rgb]{0.68,0.00,0.00}{#1}}
\newcommand{\BuiltInTok}[1]{\textcolor[rgb]{0.00,0.23,0.31}{#1}}
\newcommand{\CharTok}[1]{\textcolor[rgb]{0.13,0.47,0.30}{#1}}
\newcommand{\CommentTok}[1]{\textcolor[rgb]{0.37,0.37,0.37}{#1}}
\newcommand{\CommentVarTok}[1]{\textcolor[rgb]{0.37,0.37,0.37}{\textit{#1}}}
\newcommand{\ConstantTok}[1]{\textcolor[rgb]{0.56,0.35,0.01}{#1}}
\newcommand{\ControlFlowTok}[1]{\textcolor[rgb]{0.00,0.23,0.31}{\textbf{#1}}}
\newcommand{\DataTypeTok}[1]{\textcolor[rgb]{0.68,0.00,0.00}{#1}}
\newcommand{\DecValTok}[1]{\textcolor[rgb]{0.68,0.00,0.00}{#1}}
\newcommand{\DocumentationTok}[1]{\textcolor[rgb]{0.37,0.37,0.37}{\textit{#1}}}
\newcommand{\ErrorTok}[1]{\textcolor[rgb]{0.68,0.00,0.00}{#1}}
\newcommand{\ExtensionTok}[1]{\textcolor[rgb]{0.00,0.23,0.31}{#1}}
\newcommand{\FloatTok}[1]{\textcolor[rgb]{0.68,0.00,0.00}{#1}}
\newcommand{\FunctionTok}[1]{\textcolor[rgb]{0.28,0.35,0.67}{#1}}
\newcommand{\ImportTok}[1]{\textcolor[rgb]{0.00,0.46,0.62}{#1}}
\newcommand{\InformationTok}[1]{\textcolor[rgb]{0.37,0.37,0.37}{#1}}
\newcommand{\KeywordTok}[1]{\textcolor[rgb]{0.00,0.23,0.31}{\textbf{#1}}}
\newcommand{\NormalTok}[1]{\textcolor[rgb]{0.00,0.23,0.31}{#1}}
\newcommand{\OperatorTok}[1]{\textcolor[rgb]{0.37,0.37,0.37}{#1}}
\newcommand{\OtherTok}[1]{\textcolor[rgb]{0.00,0.23,0.31}{#1}}
\newcommand{\PreprocessorTok}[1]{\textcolor[rgb]{0.68,0.00,0.00}{#1}}
\newcommand{\RegionMarkerTok}[1]{\textcolor[rgb]{0.00,0.23,0.31}{#1}}
\newcommand{\SpecialCharTok}[1]{\textcolor[rgb]{0.37,0.37,0.37}{#1}}
\newcommand{\SpecialStringTok}[1]{\textcolor[rgb]{0.13,0.47,0.30}{#1}}
\newcommand{\StringTok}[1]{\textcolor[rgb]{0.13,0.47,0.30}{#1}}
\newcommand{\VariableTok}[1]{\textcolor[rgb]{0.07,0.07,0.07}{#1}}
\newcommand{\VerbatimStringTok}[1]{\textcolor[rgb]{0.13,0.47,0.30}{#1}}
\newcommand{\WarningTok}[1]{\textcolor[rgb]{0.37,0.37,0.37}{\textit{#1}}}

\usepackage{longtable,booktabs,array}
\newcounter{none} % for unnumbered tables
\usepackage{calc} % for calculating minipage widths
\usepackage{caption}
% Make caption package work with longtable
\makeatletter
\def\fnum@table{\tablename~\thetable}
\makeatother
\usepackage{graphicx}
\makeatletter
\newsavebox\pandoc@box
\newcommand*\pandocbounded[1]{% scales image to fit in text height/width
  \sbox\pandoc@box{#1}%
  \Gscale@div\@tempa{\textheight}{\dimexpr\ht\pandoc@box+\dp\pandoc@box\relax}%
  \Gscale@div\@tempb{\linewidth}{\wd\pandoc@box}%
  \ifdim\@tempb\p@<\@tempa\p@\let\@tempa\@tempb\fi% select the smaller of both
  \ifdim\@tempa\p@<\p@\scalebox{\@tempa}{\usebox\pandoc@box}%
  \else\usebox{\pandoc@box}%
  \fi%
}
% Set default figure placement to htbp
\def\fps@figure{htbp}
\makeatother





\setlength{\emergencystretch}{3em} % prevent overfull lines

\providecommand{\tightlist}{%
  \setlength{\itemsep}{0pt}\setlength{\parskip}{0pt}}



 


\makeatletter
\@ifpackageloaded{tcolorbox}{}{\usepackage[skins,breakable]{tcolorbox}}
\@ifpackageloaded{fontawesome5}{}{\usepackage{fontawesome5}}
\definecolor{quarto-callout-color}{HTML}{909090}
\definecolor{quarto-callout-note-color}{HTML}{0758E5}
\definecolor{quarto-callout-important-color}{HTML}{CC1914}
\definecolor{quarto-callout-warning-color}{HTML}{EB9113}
\definecolor{quarto-callout-tip-color}{HTML}{00A047}
\definecolor{quarto-callout-caution-color}{HTML}{FC5300}
\definecolor{quarto-callout-color-frame}{HTML}{acacac}
\definecolor{quarto-callout-note-color-frame}{HTML}{4582ec}
\definecolor{quarto-callout-important-color-frame}{HTML}{d9534f}
\definecolor{quarto-callout-warning-color-frame}{HTML}{f0ad4e}
\definecolor{quarto-callout-tip-color-frame}{HTML}{02b875}
\definecolor{quarto-callout-caution-color-frame}{HTML}{fd7e14}
\makeatother
\makeatletter
\@ifpackageloaded{caption}{}{\usepackage{caption}}
\AtBeginDocument{%
\ifdefined\contentsname
  \renewcommand*\contentsname{Table of contents}
\else
  \newcommand\contentsname{Table of contents}
\fi
\ifdefined\listfigurename
  \renewcommand*\listfigurename{List of Figures}
\else
  \newcommand\listfigurename{List of Figures}
\fi
\ifdefined\listtablename
  \renewcommand*\listtablename{List of Tables}
\else
  \newcommand\listtablename{List of Tables}
\fi
\ifdefined\figurename
  \renewcommand*\figurename{Figure}
\else
  \newcommand\figurename{Figure}
\fi
\ifdefined\tablename
  \renewcommand*\tablename{Table}
\else
  \newcommand\tablename{Table}
\fi
}
\@ifpackageloaded{float}{}{\usepackage{float}}
\floatstyle{ruled}
\@ifundefined{c@chapter}{\newfloat{codelisting}{h}{lop}}{\newfloat{codelisting}{h}{lop}[chapter]}
\floatname{codelisting}{Listing}
\newcommand*\listoflistings{\listof{codelisting}{List of Listings}}
\makeatother
\makeatletter
\makeatother
\makeatletter
\@ifpackageloaded{caption}{}{\usepackage{caption}}
\@ifpackageloaded{subcaption}{}{\usepackage{subcaption}}
\makeatother

\usepackage{bookmark}
\IfFileExists{xurl.sty}{\usepackage{xurl}}{} % add URL line breaks if available
\urlstyle{same}
\hypersetup{
  pdftitle={Converting Jupyter Notebooks to PDF},
  hidelinks,
  pdfcreator={LaTeX via pandoc}}


\title{\texorpdfstring{Converting Jupyter Notebooks to
PDF}{Converting Jupyter Notebooks to PDF}}
\author{}
\date{}

\begin{document}
\frame{\titlepage}


\begin{frame}[fragile]{Overview}
\protect\phantomsection\label{overview}
You've written a Jupyter notebook with code, figures, and LaTeX math ---
now you need to submit it as a PDF. This tutorial covers several ways to
do that, from one-click exports to command-line workflows.

The central challenge is that a notebook contains \textbf{code outputs}
(tables, plots) and \textbf{rich text} (Markdown, LaTeX math).
Converting both to a single PDF requires a pipeline that can handle each
piece. Some methods require LaTeX to be installed; others use your
browser's print engine instead.

\begin{tcolorbox}[enhanced jigsaw, arc=.35mm, bottomrule=.15mm, bottomtitle=1mm, breakable, colback=white, colbacktitle=quarto-callout-note-color!10!white, colframe=quarto-callout-note-color-frame, coltitle=black, left=2mm, leftrule=.75mm, opacityback=0, opacitybacktitle=0.6, rightrule=.15mm, title=\textcolor{quarto-callout-note-color}{\faInfo}\hspace{0.5em}{Do you need LaTeX installed?}, titlerule=0mm, toprule=.15mm, toptitle=1mm]

{\def\LTcaptype{none} % do not increment counter
\begin{longtable}[]{@{}lc@{}}
\toprule\noalign{}
Method & LaTeX required? \\
\midrule\noalign{}
\endhead
VS Code + Jupyter extension → PDF & Yes \\
\texttt{nbconvert\ -\/-to\ pdf} & Yes \\
\texttt{nbconvert\ -\/-to\ webpdf} & No \\
HTML → browser Print to PDF & No \\
Google Colab → Print & No \\
Quarto \texttt{-\/-to\ pdf} & Yes \\
\bottomrule\noalign{}
\end{longtable}
}

If you don't want to install LaTeX (\textasciitilde4 GB for TeX Live),
use one of the \textbf{No} options.

\end{tcolorbox}
\end{frame}

\begin{frame}[fragile]{Option 1: VS Code with the Jupyter Extension}
\protect\phantomsection\label{option-1-vs-code-with-the-jupyter-extension}
This is the most convenient if you already use VS Code for notebooks.

\begin{block}{Requirements}
\protect\phantomsection\label{requirements}
\begin{enumerate}
\tightlist
\item
  \textbf{VS Code} with the
  \href{https://marketplace.visualstudio.com/items?itemName=ms-toolsai.jupyter}{Jupyter
  extension} (\texttt{ms-toolsai.jupyter})
\item
  \textbf{Python} with \texttt{jupyter} and \texttt{nbconvert} installed
  (\texttt{pip\ install\ jupyter\ nbconvert})
\item
  \textbf{LaTeX} --- install one of:

  \begin{itemize}
  \tightlist
  \item
    \textbf{macOS}: \href{https://tug.org/mactex/}{MacTeX}
    (\textasciitilde4 GB) or the smaller
    \href{https://tug.org/mactex/morepackages.html}{BasicTeX}
    (\textasciitilde90 MB)
  \item
    \textbf{Windows}: \href{https://miktex.org/}{MiKTeX}
  \item
    \textbf{Linux}:
    \texttt{sudo\ apt\ install\ texlive-xetex\ texlive-latex-extra}
  \end{itemize}
\end{enumerate}
\end{block}

\begin{block}{Steps}
\protect\phantomsection\label{steps}
\begin{enumerate}
\tightlist
\item
  Open your \texttt{.ipynb} file in VS Code.
\item
  Click the \textbf{\ldots{}} (more actions) menu in the top-right of
  the notebook toolbar.
\item
  Select \textbf{Export} → \textbf{PDF}.
\item
  If this is your first time, VS Code will prompt you to install
  \texttt{nbconvert} if it's missing. Let it do so.
\item
  The PDF will be saved next to your notebook.
\end{enumerate}
\end{block}

\begin{block}{Troubleshooting}
\protect\phantomsection\label{troubleshooting}
\begin{itemize}
\tightlist
\item
  \textbf{``LaTeX not found''}: Make sure \texttt{pdflatex} or
  \texttt{xelatex} is on your PATH. Try \texttt{which\ pdflatex} in a
  terminal.
\item
  \textbf{Missing LaTeX packages}: The conversion may fail with errors
  about missing \texttt{.sty} files. Install the missing package through
  your LaTeX distribution's package manager
  (\texttt{tlmgr\ install\ \textless{}package\textgreater{}} for TeX
  Live).
\item
  \textbf{Plot rendering issues}: If your plots use special fonts or
  backends (e.g., PythonPlot in Julia), try re-rendering them with a
  vector-friendly backend before export.
\end{itemize}
\end{block}
\end{frame}

\begin{frame}[fragile]{Option 2: \texttt{nbconvert} from the Command
Line}
\protect\phantomsection\label{option-2-nbconvert-from-the-command-line}
If you prefer the terminal or want to automate conversion, use
\texttt{nbconvert}.

\begin{block}{Requirements}
\protect\phantomsection\label{requirements-1}
\begin{itemize}
\tightlist
\item
  Python with \texttt{jupyter} and \texttt{nbconvert}:
  \texttt{pip\ install\ jupyter\ nbconvert}
\item
  LaTeX (same as Option 1)
\end{itemize}
\end{block}

\begin{block}{Steps}
\protect\phantomsection\label{steps-1}
\begin{Shaded}
\begin{Highlighting}[]
\ExtensionTok{jupyter}\NormalTok{ nbconvert }\AttributeTok{{-}{-}to}\NormalTok{ pdf notebook.ipynb}
\end{Highlighting}
\end{Shaded}

This creates \texttt{notebook.pdf} in the same directory. To keep
intermediate files for debugging:

\begin{Shaded}
\begin{Highlighting}[]
\ExtensionTok{jupyter}\NormalTok{ nbconvert }\AttributeTok{{-}{-}to}\NormalTok{ pdf notebook.ipynb }\AttributeTok{{-}{-}debug}
\end{Highlighting}
\end{Shaded}
\end{block}

\begin{block}{Customizing the output}
\protect\phantomsection\label{customizing-the-output}
You can pass a template to control the appearance:

\begin{Shaded}
\begin{Highlighting}[]
\ExtensionTok{jupyter}\NormalTok{ nbconvert }\AttributeTok{{-}{-}to}\NormalTok{ pdf }\AttributeTok{{-}{-}template}\NormalTok{ classic notebook.ipynb}
\end{Highlighting}
\end{Shaded}

To hide code cells and show only outputs (useful for reports):

\begin{Shaded}
\begin{Highlighting}[]
\ExtensionTok{jupyter}\NormalTok{ nbconvert }\AttributeTok{{-}{-}to}\NormalTok{ pdf }\AttributeTok{{-}{-}no{-}input}\NormalTok{ notebook.ipynb}
\end{Highlighting}
\end{Shaded}
\end{block}
\end{frame}

\begin{frame}[fragile]{Option 3: \texttt{nbconvert\ -\/-to\ webpdf} (No
LaTeX Needed)}
\protect\phantomsection\label{option-3-nbconvert---to-webpdf-no-latex-needed}
This option uses a headless Chromium browser to render the notebook as a
PDF, bypassing LaTeX entirely.

\begin{block}{Requirements}
\protect\phantomsection\label{requirements-2}
\begin{itemize}
\tightlist
\item
  Python with \texttt{jupyter} and \texttt{nbconvert}:
  \texttt{pip\ install\ jupyter\ nbconvert}
\item
  \href{https://playwright.dev/}{Playwright} for Chromium:
  \texttt{pip\ install\ playwright} then
  \texttt{playwright\ install\ chromium}
\end{itemize}
\end{block}

\begin{block}{Steps}
\protect\phantomsection\label{steps-2}
\begin{Shaded}
\begin{Highlighting}[]
\ExtensionTok{jupyter}\NormalTok{ nbconvert }\AttributeTok{{-}{-}to}\NormalTok{ webpdf notebook.ipynb}
\end{Highlighting}
\end{Shaded}

The output PDF will closely match what you see when you run
\texttt{jupyter\ notebook} in your browser. This is a good option if you
have complex HTML/CSS in your notebook that LaTeX would struggle with.
\end{block}

\begin{block}{Limitations}
\protect\phantomsection\label{limitations}
\begin{itemize}
\tightlist
\item
  Page breaks may be less clean than LaTeX-produced PDFs
\item
  Very large notebooks can be slow to render
\item
  No native support for LaTeX-specific features like equation numbering
  or cross-references
\end{itemize}
\end{block}
\end{frame}

\begin{frame}[fragile]{Option 4: Convert to HTML, Then Print to PDF}
\protect\phantomsection\label{option-4-convert-to-html-then-print-to-pdf}
This is the \textbf{zero-install} method --- you only need a browser.

\begin{block}{Steps}
\protect\phantomsection\label{steps-3}
\begin{enumerate}
\item
  \textbf{Export to HTML} from VS Code (\textbf{\ldots{}} →
  \textbf{Export} → \textbf{HTML}) or from the command line:

\begin{Shaded}
\begin{Highlighting}[]
\ExtensionTok{jupyter}\NormalTok{ nbconvert }\AttributeTok{{-}{-}to}\NormalTok{ html notebook.ipynb}
\end{Highlighting}
\end{Shaded}
\item
  \textbf{Open the \texttt{.html} file} in Chrome, Firefox, or Safari.
\item
  \textbf{Print to PDF}:

  \begin{itemize}
  \tightlist
  \item
    Chrome/Edge: \texttt{Ctrl+P} (Windows) / \texttt{Cmd+P} (macOS) →
    Destination: \textbf{Save as PDF}
  \item
    Firefox: \texttt{Ctrl+P} / \texttt{Cmd+P} → \textbf{Save to PDF}
  \item
    Safari: \texttt{Cmd+P} → \textbf{PDF} dropdown → \textbf{Save as
    PDF}
  \end{itemize}
\item
  In the print dialog, enable \textbf{Background Graphics} (Chrome) or
  \textbf{Print Backgrounds} (Firefox) so code cell formatting is
  preserved.
\end{enumerate}
\end{block}

\begin{block}{Tips}
\protect\phantomsection\label{tips}
\begin{itemize}
\tightlist
\item
  Adjust \textbf{Margins} to \textbf{None} or \textbf{Minimum} for a
  cleaner look.
\item
  If the notebook is wide, try \textbf{Landscape} orientation.
\item
  You can add a \texttt{@media\ print} CSS rule to a custom template to
  control exactly how the printed PDF looks.
\end{itemize}
\end{block}
\end{frame}

\begin{frame}[fragile]{Option 5: Google Colab and Other Web-Based Tools}
\protect\phantomsection\label{option-5-google-colab-and-other-web-based-tools}
\begin{block}{Google Colab}
\protect\phantomsection\label{google-colab}
If your notebook is in Google Drive or you upload it to
\href{https://colab.research.google.com/}{Google Colab}:

\begin{enumerate}
\tightlist
\item
  Open the notebook in Colab.
\item
  \textbf{File} → \textbf{Print} (or \texttt{Ctrl+P} / \texttt{Cmd+P}).
\item
  In the print dialog, choose \textbf{Save as PDF}.
\item
  Enable \textbf{Background Graphics}.
\end{enumerate}

Colab also has \textbf{File} → \textbf{Download} → \textbf{PDF via
LaTeX} (requires a Colab Pro subscription) or \textbf{Download as
\texttt{.ipynb}} (to use one of the other methods).
\end{block}

\begin{block}{GitHub + nbviewer}
\protect\phantomsection\label{github-nbviewer}
If your notebook is on GitHub:

\begin{enumerate}
\tightlist
\item
  Copy the notebook's GitHub URL.
\item
  Paste it into \href{https://nbviewer.org}{nbviewer.org}.
\item
  Click the \textbf{Download} icon and select \textbf{PDF via LaTeX} (if
  LaTeX is available on nbviewer's backend) or download as HTML and
  print.
\end{enumerate}
\end{block}

\begin{block}{Deepnote}
\protect\phantomsection\label{deepnote}
\href{https://deepnote.com}{Deepnote} offers a built-in PDF export under
\textbf{File} → \textbf{Export as PDF}.
\end{block}
\end{frame}

\begin{frame}[fragile]{Option 6: Quarto}
\protect\phantomsection\label{option-6-quarto}
If you have \href{https://quarto.org/}{Quarto} installed, it provides
the most control over PDF output:

\begin{Shaded}
\begin{Highlighting}[]
\ExtensionTok{quarto}\NormalTok{ render notebook.ipynb }\AttributeTok{{-}{-}to}\NormalTok{ pdf}
\end{Highlighting}
\end{Shaded}

Quarto uses LaTeX (via \texttt{xelatex} or \texttt{pdflatex}) by default
but gives you far more control over formatting. You can add a YAML
header to the notebook to customize margins, fonts, and include
headers/footers --- similar to what this tutorial series uses for
homework assignments.

\begin{block}{Without LaTeX (Quarto)}
\protect\phantomsection\label{without-latex-quarto}
Quarto can also produce a PDF via Typst (a modern LaTeX alternative,
\textasciitilde30 MB):

\begin{Shaded}
\begin{Highlighting}[]
\ExtensionTok{quarto}\NormalTok{ render notebook.ipynb }\AttributeTok{{-}{-}to}\NormalTok{ typst}
\end{Highlighting}
\end{Shaded}

Or via WebTeX (requires Chrome/Chromium):

\begin{Shaded}
\begin{Highlighting}[]
\ExtensionTok{quarto}\NormalTok{ render notebook.ipynb }\AttributeTok{{-}{-}to}\NormalTok{ pdf }\AttributeTok{{-}{-}pdf{-}engine}\NormalTok{ wkhtmltopdf}
\end{Highlighting}
\end{Shaded}
\end{block}
\end{frame}

\begin{frame}{Quick Decision Guide}
\protect\phantomsection\label{quick-decision-guide}
{\def\LTcaptype{none} % do not increment counter
\begin{longtable}[]{@{}ll@{}}
\toprule\noalign{}
You want\ldots{} & Use \\
\midrule\noalign{}
\endhead
One click in VS Code & Option 1 \\
No LaTeX install, quick result & Option 3 or 4 \\
Automation / scripts & Option 2 or 6 \\
Zero local tools (browser only) & Option 5 (Colab) \\
Full control over PDF formatting & Option 6 (Quarto) \\
Submission-ready with headers/footers & Option 6 (Quarto) \\
\bottomrule\noalign{}
\end{longtable}
}
\end{frame}




\end{document}
